\section{Documentación y uso de los resultados}
\label{sec:documentacion}
% Origen: estructura del repositorio y scripts del TPI (2026-11-10).

Todo el estudio está en un repositorio Git: el modelo, los datos procesados, los scripts y este informe. La
tabla~\ref{tab:archivos} resume qué hace cada pieza.

\begin{table}[H]
\centering
\caption{Archivos del estudio}
\label{tab:archivos}
\small
\begin{tabularx}{\textwidth}{@{}>{\raggedright\arraybackslash}p{4.6cm} X@{}}
\toprule
Archivo & Función \\
\midrule
\texttt{SubteConstitucion.alp} & Modelo de AnyLogic: capa lógica, capa peatonal y experimentos \\
\texttt{verificar\_modelo.py} & Verificador automático: estructura, datos embebidos, geometría y funciones compiladas \\
\texttt{datos/embeber\_dias\_sbase.py} & Embebe en el modelo los 75 días hábiles y sus llegadas reales de trenes \\
\texttt{cargar\_corridas.py} & Valida las filas de las corridas y las carga en la planilla \\
\texttt{04-resultados-corridas.xlsx} & Planilla con las corridas e intervalos de confianza por comparación \\
\texttt{analisis\_corridas.py} & Diferencias apareadas, Bonferroni y $n^{*}$, independiente de la planilla \\
\texttt{validacion\_corridas.py} & Caudal por ventana contra SBASE y reparto entre molinetes \\
\texttt{resultados\_informe.py} & Genera las tablas y figuras de resultados de este informe \\
\texttt{informe/} & Fuente \LaTeX{} del informe y script para compilarlo en Overleaf \\
\bottomrule
\end{tabularx}
\end{table}

Para reproducir los resultados:
\begin{enumerate}
  \item correr \texttt{verificar\_modelo.py} y comprobar que todos los controles den OK;
  \item abrir el modelo en AnyLogic~8.9 PLE y ejecutar \texttt{PeatonalCorridasApareadas}, que escribe
        \texttt{corridas\_peatonales.csv} y \texttt{corridas\_peatonales\_ventanas.csv};
  \item cargar las corridas con \texttt{cargar\_corridas.py corridas\_peatonales.csv -{}-escribir};
  \item correr \texttt{analisis\_corridas.py}, \texttt{validacion\_corridas.py} y \texttt{resultados\_informe.py}.
\end{enumerate}
Con las mismas semillas, la demanda es idéntica en cualquier máquina; el movimiento de los peatones puede
diferir en décimas de segundo (sección~\ref{sec:paso6-validacion}).

\pendiente{agregar el enlace al video y, si se publica, al modelo en AnyLogic Cloud.}
